\begin{remark}[the exact identifications, with their published verification loci]\label{rem:exact-identification-checklist}
The forced-collision theorem and the sheet-resolution factorization are
stated on objects that this paper \emph{identifies} with objects of
the Kahn/KL sources.  Six identifications carry the entire
weight of the word ``canonical''.  Three are now proved in the body
(Lemma~\ref{lem:local-pair},
Definition~\ref{def:shared-neighborhood},
Lemma~\ref{lem:completion-injective}); the remaining four are listed
here with the published locus at which each was checked and the
outcome of that check.
\begin{enumerate}
\item \textbf{\(H_s/H_t\) versus Kahn's covering subgroups.}
\emph{Status: confirmed.}
\emph{Paper-side objects:} the based subgroups
\(H_s=\operatorname{im}[\pi_1(D_A\sm A,z_*)\to\pi_1(\Omega_s,z_*)]\)
and the target analogue (Definition~\ref{def:hs-ht}).
\emph{Published locus:} Kahn~I, \S7.4.1, Theorem~7.12 and its proof,
p.~44 --- the localization cover of \(S\sm A\) is the cover
corresponding to the fundamental group of \(U\sm A\) for \(U\) a local
neighborhood of \(A\); the general (possibly disconnected) case is on
the same page; the dynamical instance is \S7.4.2, p.~45.
\emph{Checked:} (i) Kahn specifies the cover by an \emph{unbased}
subgroup (a conjugacy class); the based \(H_s,H_t\) are representatives
fixed in Definition~\ref{def:shared-neighborhood}, and the collision
certificate is invariant under change of representative because the
winding homomorphism kills every conjugate of \(H_s\); (ii)
independence of the choice of local neighborhood follows from the
deformation-retraction property of local pairs (Kahn~I, \S7.4.1,
p.~44); (iii) these subgroups are \emph{not} the explicitly described
``loops that never cross \(\gamma_p\)'' subgroup of Kahn's
\emph{other} (satellite-compatible) renormalization, Kahn~I, \S7.4.4,
p.~46, which is a different construction with a different removal set.
Note also that Kahn requires of a local neighborhood only that every
component of its complement of \(A\) be an annulus; the Jordan choice
made in this paper is our own specialization, justified by
Lemma~\ref{lem:local-pair}.
\item \textbf{Excess-component position in the canonical seed.}
\emph{Status: confirmed.}
\emph{Paper-side objects:} the component \(X\ne A\) of \(F^{-1}A\)
with \(X\cap B=\varnothing\), and the removal asymmetry
\(X\subset F^{-1}B\), \(X\cap B=\varnothing\)
(Proposition~\ref{prop:excess}, Theorem~\ref{thm:excess-count}).
\emph{Published locus:} nesting and properness of the seed, Kahn~I,
\S6, p.~31 (\(U^r=f^{-r}(U^0)\), each \(f\colon U^{r+1}\to U^r\)
quadratic-like, so \(F=f^p\) is proper of degree \(2^p\)); the little
filled Julia sets \(K_p\), \(K_p(j)=f^j(K_p)\), \(\KK_p\), and
primitivity as pairwise disjointness, Kahn~I, \S1.1; the canonical
choice \(A=K_p\), \(B=\KK_p\) for \((i^p,f^p)\) with central local
degree two, Kahn~I, \S7.4.3, p.~45; and the removal asymmetry, which
is literally Kahn's \(T=(\BU\sm B)\cup A\) against
\(T'=(\BU'\sm f^{-1}B)\cup A'\), Kahn~I, \S7.4.2, p.~45.
\emph{Checked:} the central component has local degree two, so
\(2^p-2>0\) forces an excess component for \(p>1\); the invariance
\(F(K_p(j))=K_p(j)\) plus primitive disjointness gives
\(X\cap B=\varnothing\); and the source-side removal
\(X\subset F^{-1}B\) with target-side non-removal is exactly the
\(T'/T\) mismatch.  The construction meant throughout is the one of
\S7.4.3, not either variant of \S7.4.4, p.~46
(Definition~\ref{def:canonical}).
\item \textbf{Collision surviving every equalizer.}
\emph{Status: confirmed.}
\emph{Paper-side step:} inductive propagation of the collision pair
through the pullback tower in equalizer coordinates
(Theorem~\ref{thm:propagation}).
\emph{Published locus:} Kahn~I, \S7.3.2, pp.~42--43: the equalizer
coordinates \(\BU''=\{(x,y):\ i(x)=f(y)\}\), \(i_1(x,y)=y\),
\(f_1(x,y)=x\); the iff-chain giving
\(K''=i_1^{-1}(K')=f_1^{-1}(K')\); all critical points of \(f_1\) in
\(K''\); the properness statement ``\(f_1\) is proper of degree
\(d\)''; and the commuting tower \(i_n\circ f_{n+1}=f_n\circ i_{n+1}\)
with its iterates \(i^n\).  That each stage is again a degree-\(d\)
\(\psi\)-polynomial-like map is Lemma~7.11, p.~42.
\emph{Checked:} \(f_n\) restricted to the complements of the filled
sets is an unbranched covering (branch values in \(K^n\) by
Definition~7.1 and the remark after it, p.~36, and
\(f_n^{-1}(K^n)=K^{n+1}\) by Definition~7.1(4)); it is \emph{onto}
each complementary annulus by properness of positive degree
(Definition~7.1(3), p.~36; Lemma~7.9, p.~40; \S7.3.2, p.~42) together
with connectedness of \(A_n\); the chosen pairs \((x_n^\pm,y_n)\)
satisfy the equalizer membership condition \(i_n(x_n^\pm)=f_n(y_n)\);
the leg assignment is Kahn's, and reversing it would break the
induction; and the propagated points remain in the complementary
annuli at every depth.  The immersion hypothesis needed for
Lemma~7.8, p.~40 (without which, as the source's own example on p.~41
shows, the set-theoretic equalizer need not be the pullback in the
category of complex \(1\)-manifolds) is regenerated at each stage by
Proposition~\ref{prop:base-change}(6) and independently by Kahn~I,
Lemma~7.7, p.~40 and Lemma~7.9, p.~40.
\item \textbf{Direct-decoration depths and tile transport.}
\emph{Status: confirmed} (the transport itself remains conditional on
strong stage alignment; see below).
\emph{Paper-side step:} noninjectivity of the ambient annulus maps
\(I_{m,n}\), of their universal-cover tile lifts, and of the glued
pseudo-piece immersions at the depths \(m=\nu q+1\),
\(n=(\nu-1)q+1\) (Theorem~\ref{thm:propagation}).
\emph{Published locus:} KL~II, \S2.1, p.~6 (\(P=Y^{(\nu-1)q+1}\); the
two univalent \(f^q\)-pullbacks \(Q_L,Q_R\subset P\) of depth
\(q\nu+1\)) and Figure~2.1, p.~7; KL~II, \S3.2.1, pp.~11--12
(pseudo-piece as path classes, and equivalently as \(Y\) with copies
of the universal cover of \(\BU^m\sm K(f)\) glued on);
Lemma~3.3, p.~12 (the glued copies embed, with pairwise disjoint
images); Lemma~3.4, pp.~12--13 (the inclusion-induced arrow is an
immersion, with component index map ``surjective but not necessarily
injective''; the dynamical arrow is a branched covering).
\emph{Checked:} the paper's ambient map at depths \((m,n)\) is the
\emph{inclusion-induced} arrow of KL~II, not the dynamical arrow (which
is an isomorphism on univalent pieces and is not contested); the depth
arithmetic matches \S2.1 exactly; the tile objects are universal covers
of exactly the annuli \(A_m\); the collision transported from the
canonical seed lands in the correct glued copy, using the
disjoint-image property of Lemma~3.3, p.~12 --- not Lemma~3.4, which
supplies only immersions; and the realization quantifier is kept
separate.  Two scope facts belong here: KL~II never forms a
pseudo-piece from \(Q_L\) or \(Q_R\)
(Remark~\ref{rem:QLQR-scope}), and the annular degree \(+1\) used in
the transport is \emph{derived} in Lemma~\ref{lem:annular-degree} from
Kahn~I, Theorem~7.3, p.~37, Lemma~7.4(2), p.~37 and KL~II, Lemma~3.1,
p.~10, and is stated outright for the canonical renormalization in the
proof of Kahn~I, Theorem~8.11, p.~51.
\end{enumerate}
The residual open item is not on this list: it is the
seed-identification reading of the strong stage-alignment hypothesis,
discussed in Remark~\ref{rem:alignment-is-seed-identification}.
\end{remark}

\begin{proof}[Source verification (documentary; not a mathematical proof)]
Every locus above was checked against the published texts: Kahn~I in
arXiv:math/0609045v3 (28 December 2023, 63 pp.), whose sectioning is
that used throughout this paper, and KL~II in
Ann.\ Sci.\ \'Ec.\ Norm.\ Sup\'er.\ (4) \textbf{41} (2008) 57--84 and
the identically sectioned arXiv:math/0609046v2 (12 February 2007).
The SHA-256--locked local \TeX{} archives from which the paraphrases of
Section~\ref{sec:background} were originally prepared correspond to
exactly those two versions; their line-number concordance is recorded
in Appendix~\ref{app:sourcelock}.  The authoritative statements of the
forced-collision theorem and the reflection criterion are
Theorems~\ref{thm:forced-collision}
and~\ref{thm:subgroup-reflection-criterion} as proved in
Sections~\ref{sec:criteria}--\ref{sec:resolution} of this paper.  This
checklist adds no new mathematical claim.
\end{proof}
